% correlatividades geografia
\thispagestyle{empty}
\PageDecor{\CorrPageNo}
\SectionTag{CORRELATIVIDADES}
\DocHeading{Profesorado de Educación Secundaria en Geografía}

\begingroup
\small
\setlength{\tabcolsep}{4.5pt}
\renewcommand{\arraystretch}{1.34}
\setlength{\LTleft}{0pt}
\setlength{\LTright}{0pt}
\setlength{\LTpre}{0.35em}
\setlength{\LTpost}{0.75em}

\newcommand{\GeoTableStart}{
\begin{longtable}{|>{\RaggedRight\arraybackslash}p{0.30\textwidth}|>{\RaggedRight\arraybackslash}p{0.64\textwidth}|}
\hline
\textbf{Unidad Curricular} & \textbf{Correlativa con} \\
\hline
\endfirsthead
\hline
\textbf{Unidad Curricular} & \textbf{Correlativa con} \\
\hline
\endhead
}

\newcommand{\GeoTableEnd}{
\end{longtable}
}

\newcommand{\GeoRow}[2]{
\rule{0pt}{2.55ex}#1 & #2 \\
\hline
}

\DocQuestion{Tabla 1. Correlatividades de 2\textsuperscript{o} año}
\GeoTableStart
\GeoRow{Filosofía}{Pedagogía (R).}
\GeoRow{Historia Social y Política Argentina y Latinoamericana}{Pedagogía (R).}
\GeoRow{Psicología Educacional}{Pedagogía (R).}
\GeoRow{Práctica Docente II}{Práctica Docente I (A); Pedagogía (R); Didáctica General (R).}
\GeoRow{Sujetos de la Educación Secundaria}{Didáctica General (R); Pedagogía (R).}
\GeoRow{Didáctica de las Ciencias Sociales}{Didáctica General (A); Espacios Urbanos y Rurales del Mundo Contemporáneo (R).}
\GeoRow{Organización del Espacio Geográfico Americano}{Paisajes Geográficos Mundiales; Naturaleza y Sociedad I (A).}
\GeoRow{Sistema Urbano y Desarrollo Rural Argentino}{Espacios Urbanos y Rurales del Mundo Contemporáneo (R).}
\GeoRow{Turismo y Desarrollo Local}{Geografía de Entre Ríos (R).}
\GeoTableEnd

\DocQuestion{Tabla 2. Correlatividades de 3\textsuperscript{o} año}
\GeoTableStart
\GeoRow{Historia de la Educación Argentina}{Pedagogía (A); Historia Social y Política Argentina y Latinoamericana (R).}
\GeoRow{Sociología de la Educación}{Historia Social y Política Argentina y Latinoamericana (R).}
\GeoRow{Práctica Docente III}{Práctica Docente II (A); Psicología Educacional (R); Sujetos de la Educación Secundaria (R); Didáctica de las Ciencias Sociales (R).}
\GeoRow{Análisis y Organización de las Instituciones Educativas}{Sujetos de la Educación Secundaria (R); Práctica I (A); Práctica Docente II (R).}
\GeoRow{Organización del Espacio Geográfico Argentino}{Organización del Espacio Americano (A); Sistema Urbano y Desarrollo Rural Argentino (R).}
\GeoRow{Naturaleza y Sociedad II}{Naturaleza y Sociedad I (A).}
\GeoRow{Epistemología de la Geografía}{Filosofía (R); Espacios Urbanos y Rurales del Mundo Contemporáneo (A).}
\GeoRow{Didáctica de la Geografía}{Didáctica de las Ciencias Sociales (A); Sujetos de la Educación Secundaria (R).}
\GeoTableEnd

\DocQuestion{Tabla 3. Correlatividades de 4\textsuperscript{o} año}
\GeoTableStart
\GeoRow{Derechos Humanos: Ética y Ciudadanía}{Hist. Social y Política Argentina y Latinoamericana (A); Filosofía (R); Sujetos de la Educación Secundaria (A).}
\GeoRow{Práctica Docente IV}{Todas las UC de Tercer año (A).}
\GeoRow{Geotecnología}{Técnicas de la Representación Cartográfica (A).}
\GeoRow{Geografía Económica}{Epistemología de la Geografía (R); Sistema Urbano y Desarrollo Rural Argentino (A); Organización del Espacio Americano (A).}
\GeoRow{Geografía Política y Cultural}{Epistemología de la geografía (R); Organización del Espacio Argentino (R).}
\GeoRow{Culturas y Problemáticas Ambientales}{Organización del Espacio Geográfico Argentino (R); Naturaleza y Sociedad II (R).}
\GeoRow{Seminario en Geografía}{Didáctica de las Ciencias Sociales (A); Didáctica de la Geografía (R); Práctica Docente III (A).}
\GeoTableEnd

\DocPara{\textbf{Referencias:} A: aprobada - R: regular.}
\DocPara{\textsuperscript{36} Para cursar una unidad curricular, el estudiante puede tener regularizada/s y/o aprobada/s las UC correlativa/s anterior/es. Para rendir una UC, el estudiante debe tener aprobada/s la/s unidades curriculares correlativas.}
\endgroup
